% ROUND7-REFEREE-POSITIVE-CLOSURE
\section{Round-seven reconstruction: coarse histories, suspension kernels, and a Riesz--Schur memory dilation}
\label{sec:r7-a4}

This section replaces the round-six controlling module.  The full microscopic
Sinai state is never used as the conditioning variable for a nondegenerate
future law.  Conditional kernels live on a stable-leaf quotient which forgets
future-determining coordinates.  Continuous-time kernels are constructed
before Ray theory.  Transmission zeros are treated by finite-rank Riesz
residues of the compressed resolvent, not by a rational Smith--McMillan claim
for a general meromorphic matrix.

\subsection{Full and coarse histories}

Let $E_{\rm mic}$ be the graph-completed suspension state and let
$\Phi_t^\varepsilon$ be its deterministic invertible flow.  Let
\[
 \pi_s:E_{\rm mic}\longrightarrow E_s
\]
be the quotient by local stable leaves together with the present slow
coordinate and the suspension clock.  Conditional measures on stable fibers
are the normalized Gibbs disintegrations $\lambda_y^\varepsilon$.  The coarse
history space is
\[
 \mathsf H_s=D(( -\infty,0],E_s),
 \qquad
 \mathcal F_t^s=\sigma\{\pi_s(\Phi_r^\varepsilon x):r\le t\}.
\]
The full microscopic history still determines a Dirac future; no theorem below
claims otherwise.

For a coarse history $h$ ending at $y=h(0)$ define
\[
 K_t^\varepsilon(h,A)=
 \int_{\pi_s^{-1}(y)}
  1_A\bigl(\pi_s(\Phi_{[0,t]}^\varepsilon x)\bigr)
  \,d\lambda_{h}^{\varepsilon}(x),
\]
where $\lambda_h^\varepsilon$ is the stable-fiber conditional law obtained by
transporting the quotient $g$-kernel along the past.  Exponential memory loss
makes this version independent of remote representatives of $h$.

\begin{lemma}[Nondegenerate coarse disintegration]
\label{lem:r7-a4-coarse}
For every regular coarse history, $\lambda_h^\varepsilon$ is a probability on
a positive-dimensional stable fiber.  The kernel $K_t^\varepsilon$ is a Dirac
mass only when the quotient fiber itself is a singleton.  Given the complete
microscopic history, the conditional kernel is explicitly the Dirac mass at
the deterministic future, and that kernel is not used in the homogenization
statement.
\end{lemma}

\begin{proof}
The local product structure disintegrates the prepared Gibbs measure into
unstable quotient and stable conditional measures with Holder Jacobians.  The
coarse past fixes the quotient itinerary but not the stable coordinate.  The
usual contraction of stable holonomies makes the conditional density a
Cauchy limit as the remote past is extended.  Conditioning on the full state
instead fixes the stable coordinate, yielding the stated Dirac law.
\end{proof}

\subsection{A continuous-time suspension semigroup}

Let $P$ be the one-step quotient $g$-kernel, $\tau$ the roof, and
$N_t(x)$ the number of returns completed by time $t$.  For bounded uniformly
continuous $F$ on $E_s$, define for every real $t\ge0$
\[
 (\mathcal P_tF)(y,u)=
 \sum_{n\ge0}\int
 F\bigl(T^nx,u+t-S_n\tau(x)\bigr)
 1_{\{S_n\tau(x)\le u+t<S_{n+1}\tau(x)\}}
 \,dP_y^{(n)}(x).
\]
This renewal formula defines irrational and rational times simultaneously.

\begin{theorem}[Coarse-history Feller semigroup]
\label{thm:r7-a4-feller}
The operators $\mathcal P_t$ form a conservative strongly continuous Feller
semigroup on the weighted uniformly continuous history state.  Their kernels
satisfy Chapman--Kolmogorov for all real $s,t\ge0$.  The resolvent
\[
 R_\lambda F=\int_0^\infty e^{-\lambda t}\mathcal P_tF\,dt
\]
is therefore a genuine sub-Markov resolvent satisfying the resolvent identity
before any Ray compactification is invoked.
\end{theorem}

\begin{proof}
Partitioning according to $N_t$ gives an absolutely convergent sum because the
roof is bounded below and has an exponential tail.  Splitting a path at time
$s$ and then at time $s+t$ proves Chapman--Kolmogorov term by term.  Stable
holonomy continuity and the roof tail give the Feller modulus; dominated
convergence proves strong continuity at zero.  Fubini and the semigroup
identity yield the resolvent identity.
\end{proof}

Choose a countable Ray cone generated by $R_\lambda F$ for a separating family
of $F$ and rational $\lambda>0$.

\begin{corollary}[Noncircular Ray realization]
\label{cor:r7-a4-ray}
The Ray compactification of the resolvent in
Theorem~\ref{thm:r7-a4-feller} carries a right-continuous strong Markov process
whose transition semigroup extends $\mathcal P_t$.  The construction does not
use the desired right continuity to define the resolvent.
\end{corollary}

\begin{proof}
The already constructed resolvent is conservative, separates points on the Ray
cone, and is normal by monotone convergence.  The standard Ray theorem applies.
\end{proof}

\subsection{History eigenfunctions and the Doob transform}

Let $L_\Psi$ be the quotient transfer operator for a bounded source $\Psi$.
On each positive recurrent component let $h_\Psi$ and $\ell_\Psi$ be its
strictly positive eigenfunction and eigenfunctional, normalized by
$\ell_\Psi(h_\Psi)=1$.  Lift $h_\Psi$ to the suspension by the renewal formula
rather than by a possibly nonexistent last return:
\[
 H_\Psi(y,u)=e^{-\lambda_\Psi u}
 \mathbb E_y\left[
 e^{\int_0^{\tau_1-u}\Psi(X_s)ds}
 h_\Psi(Y_1)
 \right].
\]
This formula is defined at every coding representative and has matching
oriented traces.

\begin{theorem}[Bounded history Doob transform]
\label{thm:r7-a4-doob}
On compact source charts,
$0<c\le H_\Psi\le C<\infty$.  The operators
\[
 \mathcal P_t^\Psi F
 =e^{-t\lambda_\Psi}H_\Psi^{-1}
 \mathcal T_t^\Psi(H_\Psi F)
\]
form a conservative Feller semigroup on the coarse history state and preserve
the Doob invariant probability.
\end{theorem}

\begin{proof}
The quotient cone theorem gives uniform upper and lower bounds for
$h_\Psi$ on the normalized inducing base.  The roof exponential moment and
bounded source transfer these bounds through the renewal formula.  The
eigenfunction identity proves conservation and the semigroup law.
\end{proof}

\subsection{Riesz--Schur dilation of transmission zeros}

Let $L_D$ be the generator of one fixed Doob semigroup on
$\mathcal H=L^2(\nu_D)$, let $P$ be the orthogonal projection onto a finite
resolved space $V\subset D(L_D^r)$, and define
\[
 C(z)=P(z-L_D)^{-1}P|_V.
\]
Assume the A2 resolvent estimate in a strip
$\operatorname{Re}z>-\gamma$ and let $\mathcal Z_0$ be the finite multiset of
zeros of $\det C(z)$ in the smaller strip
$\operatorname{Re}z\ge-\gamma_0$.  Around each zero take a contour
$\Gamma_j$ avoiding poles and define the finite-rank residue space generated
by the Laurent coefficients of $C(z)^{-1}$.

\begin{theorem}[Operator-compatible Riesz--Schur memory dilation]
\label{thm:r7-a4-riesz}
There exist a finite-dimensional space $Z$, matrices $A_Z,B_Z,C_Z$, and an
orthogonal projection $\widetilde P$ on $V\oplus Z$ such that the Schur
complement of
\[
 \widetilde C(z)=
 \widetilde P(z-\widetilde L)^{-1}\widetilde P
\]
with respect to $Z$ equals $C(z)$ on
$\operatorname{Re}z>0$.  The spectrum of $A_Z$ consists exactly of the
transmission zeros in $\mathcal Z_0$, with their algebraic multiplicities.
The residual compressed inverse is analytic and zero-free in
$\operatorname{Re}z\ge-\gamma_0$.
\end{theorem}

\begin{proof}
The analytic Fredholm theorem makes $C(z)^{-1}$ meromorphic in the strip.
Its principal parts at the finitely many zeros are finite-rank matrix
polynomials in $(z-z_j)^{-1}$.  Realize each principal part by its companion
Jordan block on a finite residue space and take their direct sum.  Subtracting
these principal parts leaves an analytic inverse.  The block operator
\[
 \widetilde L=
 \begin{pmatrix}PL_DP& C_Z\\ B_Z&A_Z\end{pmatrix}
\]
with the unresolved $Q$-resolvent retained in the $V$ block has, by the block
resolvent identity, Schur complement equal to the original compressed
resolvent.  The coefficients were extracted from the actual Laurent residues,
so this is an operator-compatible dilation rather than an arbitrary rational
factorization.  The remaining inverse is zero-free by construction.
\end{proof}

Let
\[
 \widehat K_{\rm res}(z)
 =C_{\rm res}(z)^{-1}-zI-A_{\rm res}.
\]

\begin{theorem}[Residual memory decay]
\label{thm:r7-a4-memory}
If $V\subset D(L_D^r)$ with
$r$ larger than the polynomial loss in the A2 high-frequency resolvent by two,
then
\[
 \|K_{\rm res}(t)\|\le Ce^{-\gamma_0t}.
\]
The exact resolved equation is the coupled Volterra system on $V\oplus Z$;
all slow auxiliary modes and the residual memory are obtained from the same
microscopic time scaling.
\end{theorem}

\begin{proof}
Theorem~\ref{thm:r7-a4-riesz} removes every zero and pole in the contour-shift
strip.  The $D(L_D^r)$ assumption gives a vertical-line bound integrable after
$r$ integrations by parts.  Shift the Bromwich contour to
$\operatorname{Re}z=-\gamma_0$ and estimate the horizontal sides using the A2
Dolgopyat loss.  The companion modes are part of the same accelerated
resolvent, so their scaling follows from the pole locations rather than being
inserted as an unscaled ODE.
\end{proof}

\subsection{Enhanced and conditional homogenization}

Assume the induced observable has exponential moments.  Then martingale--coboundary
decomposition supplies moments of every finite order for first and second
levels.  Let $X^\varepsilon$ denote the slow observable and
$\mathbb X^\varepsilon$ its canonical iterated integral.

\begin{theorem}[Enhanced coarse-history invariance principle]
\label{thm:r7-a4-rough}
For every $p>2$,
$(X^\varepsilon,\mathbb X^\varepsilon)$ converges in $p$-variation rough-path
topology to the diffusion rough path with Green--Kubo covariance and the
deterministic area anomaly.  Moreover, uniformly on compact subsets of
$\mathsf H_s$,
\[
 K_T^\varepsilon(h,\cdot)\Longrightarrow
 K_T^{\rm diff}(X_0(h),\cdot).
\]
This statement is conditional on the coarse quotient history.  Conditional on
the full microscopic history the kernel remains Dirac.
\end{theorem}

\begin{proof}
The all-order moment bounds give Kolmogorov estimates for increments and the
second level, hence tightness for every $p>2$.  Finite-dimensional convergence
follows from the martingale approximation; antisymmetric coboundary terms give
the area anomaly.  For conditional convergence, truncate the coarse past at a
Young return.  Stable-fiber coupling makes the truncation error exponentially
small uniformly on compact history sets.  The finite-past kernel is an
integral against $\lambda_h^\varepsilon$, so the same martingale estimates hold
uniformly.  Let the truncation recede after taking the homogenization limit.
No deterministic point kernel is approximated by a Brownian law.
\end{proof}

\begin{theorem}[Round-seven A4 closure]
\label{thm:r7-a4-main}
The conditional theory uses a genuine coarse history with stable-fiber
randomness, possesses a continuous-time Feller semigroup and noncircular Ray
resolvent, a globally defined history Doob eigenfunction, an
operator-compatible Riesz--Schur dilation of transmission zeros, exponential
residual memory under quantified generator regularity, and an enhanced
conditional invariance principle consistent with determinism.
\end{theorem}

\begin{proof}
Combine Lemma~\ref{lem:r7-a4-coarse},
Theorems~\ref{thm:r7-a4-feller},
\ref{thm:r7-a4-doob},
\ref{thm:r7-a4-riesz},
\ref{thm:r7-a4-memory}, and
\ref{thm:r7-a4-rough}, and
Corollary~\ref{cor:r7-a4-ray}.
\end{proof}
